% Real-world MCP corpus evaluation tables.
% Generated from reports/corpus-run-1781854580/corpus_summary.json
% and reports/corpus-run-1781854580/semantic_summary.json.

\begin{table}[t]
\centering
\caption{End-to-end coverage and failure breakdown on real-world MCP repositories.}
\label{tab:realworld-coverage}
\begin{tabular}{lrrrrrr}
\toprule
Subset & Resolved & Scanned & Success & Build & Startup & Call/Timeout \\
\midrule
All resolved repos & 91 & 35 & 38.5\% & 19 & 31 & 6 \\
Tier-1 MCP servers & 42 & 28 & 66.7\% & 4 & 10 & 0 \\
Tier-2 MCP-related repos & 49 & 7 & 14.3\% & 15 & 21 & 6 \\
\bottomrule
\end{tabular}
\end{table}

\begin{table}[t]
\centering
\caption{Semantic extraction coverage over the full MCP corpus.}
\label{tab:semantic-full-coverage}
\begin{tabular}{lrrrr}
\toprule
Source & Repos & Tools & Described & Sensitive \\
\midrule
Dynamic \texttt{tools/list} & 35 & 539 & 539 & 377 \\
Static source registration & 36 & 588 & 377 & 468 \\
Repo-level signal only & 27 & 0 & 0 & 0 \\
No semantic signal & 2 & 0 & 0 & 0 \\
\midrule
Total tool metadata & 71 & 1127 & 916 & 845 (75.0\%) \\
\bottomrule
\end{tabular}
\end{table}

\begin{table}[t]
\centering
\caption{Semantic capability distribution across extracted MCP tools.}
\label{tab:semantic-capabilities}
\begin{tabular}{lr}
\toprule
Capability & Tools \\
\midrule
Network access & 395 \\
Browser automation & 348 \\
Unknown / benign utility & 282 \\
Cloud or SaaS operation & 209 \\
Filesystem access & 204 \\
Code repository operation & 198 \\
Database access & 196 \\
Credential or secret handling & 94 \\
Shell or process execution & 62 \\
\bottomrule
\end{tabular}
\end{table}

\begin{table}[t]
\centering
\caption{Cross-validation between declared tool semantics and runtime network evidence.}
\label{tab:semantic-cross-validation}
\begin{tabular}{lrp{0.45\linewidth}}
\toprule
Relation & Repos & Interpretation \\
\midrule
Declared egress risk and observed egress & 5 & Confirmed by runtime evidence \\
Declared egress risk only & 17 & Declared capability was not exercised by this run \\
Observed egress only & 2 & Potential metadata understatement or classifier miss \\
Neither declared nor observed & 11 & No egress evidence in this run \\
\bottomrule
\end{tabular}
\end{table}

\paragraph{Real-world coverage.}
We evaluate SandboxScan on a corpus of 100 popular MCP-related repositories.
SandboxScan resolves 91 repositories and completes end-to-end dynamic scans
for 35 of them. Coverage is substantially higher on dedicated MCP servers:
28 out of 42 resolved Tier-1 servers are scanned successfully (66.7\%),
compared with 7 out of 49 Tier-2 MCP-related repositories (14.3\%).
Most incomplete scans are due to operational reproducibility rather than
the taint analysis itself: 31 repositories fail during MCP startup or
protocol initialization, and 19 fail during build.

\paragraph{Semantic tool profiling.}
We additionally run a semantic-only pass over the full 100-repository corpus.
This pass reuses protocol-level \texttt{tools/list} metadata when available
and falls back to static source-level MCP tool registrations for repositories
that did not complete dynamic execution. For repositories where no tool
registration can be recovered, we record only a weaker repository-level
semantic signal from README and repository metadata; these signals are not
counted as tools. Overall, SandboxScan recovers tool metadata for 71
repositories: 35 from dynamic \texttt{tools/list} responses and 36 from static
source registrations. These repositories expose 1,127 tools, of which 916
include descriptions and 845 (75.0\%) are classified as security-sensitive.
The most common declared capabilities are network access, browser automation,
cloud/SaaS operations, filesystem access, code repository operations, and
database access. We use this profile as attack-surface characterization and
test prioritization evidence; it is not interpreted as a vulnerability
finding.

\paragraph{Semantic cross-validation.}
For the 35 repositories that complete dynamic execution, we further compare
the declared capability of the specific tool invoked during the scan with
runtime network evidence collected by the egress monitor. This check is
intentionally conservative: it validates only network-observable behavior,
and a declared-only result does not indicate a false positive because the
single generated invocation may not exercise all declared behavior. Among
35 dynamically scanned cases, all have called-tool metadata. Twenty-two
invoked tools declare an egress-related capability, and seven executions
produce runtime network egress. Five cases are both declared and observed,
seventeen are declared-only, two are observed-only, and eleven have neither
declared nor observed egress evidence.

\paragraph{Anonymized case studies.}
We selected three dynamically executed repositories from the 100-repository
corpus to illustrate how the cross-validation evidence appears in individual
runs. Repo A exposed a code-review tool whose metadata implied GitHub/code
repository and network capability; the invoked tool returned an error containing
\texttt{https://api.github.com/graphql}, while the monitor recorded two denied
connections to \texttt{api.github.com:443}. Repo B exposed a web-search tool
classified as network/browser/database capable; even with invalid generated
arguments, execution attempted an outbound request to a third-party capture API,
which the monitor denied. Repo C exposed a configuration/desktop-control tool
whose called-tool metadata did not classify as network-capable, but the runtime
trace still showed a denied connection to an external browser-support host and
the tool result reflected local configuration and onboarding prompt text. These
examples are anonymized because they validate attack-surface evidence rather
than report repository-specific vulnerabilities.
